混入率が高いと水準$q$の種はIDのほぼ全数に達し，$\pN$は情報を失う（\secref{in200}）．そこで，棄却の閾値に上限を付けた規則（\secref{seed}：$c_i=\min(q\,i/m',\,\tau)$，$\tau{=}0.05$，$q{=}0.2$）を実行前に登録し，上限なしの規則と同じ実行の中で評価した（付録\ref{sec:app_hyp}）．

\textbf{ImageNet-1Kでは変わらない．}上限が効くのはStorey型の閾値が$\tau$を超えるSSB-hardとOpenImage-O（iNaturalistではわずか）で，SSB-hardでは種が折の\NCapSeedSsbHardStoreyRatio\%から\NCapSeedSsbHardStoreycRatio\%に減り，種に混じるIDの割合は\NCapSeedSsbHardStoreyIdfrac{}から\NCapSeedSsbHardStoreycIdfrac{}に下がる．near-OODの\FPR{}は単体で\NCapStoreycAloneNearFpr\%（上限なし\NAloneNearFpr\%；対応のある差\NCapStoreycAloneNearDfprVsStorey{} [\NCapStoreycAloneNearDfprVsStoreyLo, +\NCapStoreycAloneNearDfprVsStoreyHi]），$\times$TINSで\NCapStoreycNearFpr\%（+\NCapStoreycNearDfprVsStorey{}），far-OODとFour-OODの単体は\NCapStoreycAloneFarFpr\%，\NCapStoreycAloneFourFpr\%で，いずれも上限なしと0.25ポイント以内である．

\textbf{混入率が高いベンチマークでは害がほぼ消える．}ImageNet-200のnear-OODでは\FPR{}が\NInTwoZeroZeroMainNearFpr\%から\NInTwoZeroZeroStoreycNearFpr\%へ下がり（差\NInTwoZeroZeroStoreycNearDfprVsStorey{} [\NInTwoZeroZeroStoreycNearDfprVsStoreyLo, \NInTwoZeroZeroStoreycNearDfprVsStoreyHi]），$\pP$だけ（\NInTwoZeroZeroLpNearFpr\%）より\NInTwoZeroZeroStoreycNearDfprVsLpAbs{}ポイント低い（[\NInTwoZeroZeroStoreycNearDfprVsLpLo, \NInTwoZeroZeroStoreycNearDfprVsLpHi]）．far-OODは\NInTwoZeroZeroStoreycFarFpr\%で$\pP$だけ（\NInTwoZeroZeroLpFarFpr\%）と並ぶ．CIFAR-100ではnear-OODが\NCifarMainNearFpr\%から\NCifarStoreycNearFpr\%（$\pP$だけ\NCifarLpNearFpr\%），far-OODが\NCifarMainFarFpr\%から\NCifarStoreycFarFpr\%になり，$\pP$だけ（\NCifarLpFarFpr\%）との差は+\NCifarMainFarDfprVsLp{}から+\NCifarStoreycFarDfprVsLp{} [\NCifarStoreycFarDfprVsLpLo, \NCifarStoreycFarDfprVsLpHi]ポイントに縮む．\secref{rec}の格子の\NCaprecNCells{}条件でも，上限付きの規則が$\pP$だけより1ポイント以上悪い条件は混入率30\%以上の\NCaprecNHi{}条件で\NCaprecStoreycHiNWorseLpOne{}，$r\ge10$・混入率30\%未満の\NCaprecLoROneZeroN{}条件では上限なしとの差が最大\NCaprecLoROneZeroMaxStoreycMinusStorey{}ポイントと利得を保つ（付録\ref{sec:app_cap}）．

\textbf{位置づけ．}本稿の主構成は登録どおり上限なしのStorey型であり（\tabref{main}），上限付きは登録済みの変種として報告する．上限付きの種はCIFAR-100 far-OODの+\NCifarStoreycFarDfprVsLp{}ポイントを除いて$\pP$だけと同等以上で，ImageNet-1Kでは上限なしと変わらないから，混入率が未知の運用には上限付きを推す．
